% 第七章 系统实现与测试

\section{系统实现与测试}

\subsection{开发环境搭建}
\par 为保证系统可复现与可维护性，嵌入式端、服务器端与语音交互端分别采用独立的开发与依赖管理方式。嵌入式端统一使用PlatformIO进行工程管理与编译烧录，ESP32-S3工程使用Arduino框架并在 \texttt{platformio.ini} 中声明芯片型号、串口波特率与依赖库（TFT显示库、MQTT库、软串口等）。STM32工程同样使用PlatformIO与Arduino生态库完成I\(^2\)C传感器驱动与串口上报实现，便于快速验证多传感器采集链路。OpenMV侧采用MicroPython脚本方式开发，通过OpenMV IDE完成脚本烧录、串口日志查看与摄像头参数调试。
\par 服务器端采用Node.js + pnpm进行依赖管理与运行。服务入口为 \texttt{server.js}，启动后会读取 \texttt{server\_config.json} 以确定监听地址与端口、数据点上限、MQTT配置等，并加载历史数据文件与设置文件。前端页面作为静态资源由服务器托管，便于一体化部署。语音交互端采用Python运行，依赖在 \texttt{requirements.txt} 中管理并通过虚拟环境安装，核心模块包含实时ASR、LLM对话与工具调用、实时TTS以及音频采集播放等。
\par 配图建议如下。图7-1建议为“开发环境与工程结构示意图”，画法为一张分区示意：左侧列出嵌入式端（ESP32-S3/STM32/OpenMV）的工程目录与关键文件（如 \texttt{platformio.ini}、\texttt{src/main.cpp}、OpenMV脚本），右侧列出服务器端与语音助手端的目录与关键文件（如 \texttt{server.js}、\texttt{public/}、\texttt{voice\_assistant/}），底部用箭头标注三端的依赖管理方式（PlatformIO/pnpm/pip）以及运行入口命令。该图可以用Visio/diagrams.net绘制，也可以用PPT绘制后导出图片。
\begin{figure}[H]
\centering
\fbox{\parbox{0.92\linewidth}{\centering 图7-1占位：开发环境与工程结构示意图。\\
要求展示三端工程目录结构、关键配置文件与运行入口，并标注PlatformIO/pnpm/pip等依赖管理方式。}}
\caption{开发环境与工程结构示意图（建议配图）}
\label{fig:dev_env_structure}
\end{figure}

\subsection{硬件实现与调试}
\par 硬件实现阶段以“模块化搭建、分步验证”为策略，优先保证单模块可用，再进行系统级联调。采集侧以STM32驱动AHT20、BMP280与SGP40等I\(^2\)C传感器完成温湿度、气压与VOC数据读取，并通过USART2按行上报。网关侧以ESP32-S3连接TFT显示屏、MQ气体传感器模拟输入，并通过两路UART分别连接STM32与OpenMV，最终通过WiFi接入家庭网络。视觉侧OpenMV连接摄像头与存储介质，运行人脸采集与识别脚本，并通过UART与ESP32-S3完成命令交互与数据回传。
\par 调试过程遵循由易到难的顺序。首先对电源与连线进行基本检查，确保各模块供电电压与地线共地正确，并确认3.3V电平兼容。随后分别验证I\(^2\)C总线的可用性与地址识别，确保传感器初始化成功并能输出稳定读数。接着验证两条UART链路的收发正确性，确认STM32上报数据能被ESP32解析，OpenMV命令能被正确识别并产生响应。最后验证TFT显示刷新与WiFi连接，确保网关能在不阻塞主循环的情况下同时完成显示、上报与控制转发。
\par 配图建议如下。图7-2建议为“硬件连接与实物连线图”，可用两张图组合：一张是逻辑连接图（模块框图 + 接口连线 + 引脚号），另一张是实物照片（包含STM32板、ESP32-S3板、OpenMV板、TFT屏与传感器模块的实际连线），并在图中用标注框指出UART1/UART2、I\(^2\)C、SPI与电源线的走向。图7-3建议为“I\(^2\)C总线与传感器在线性验证截图/照片”，例如OpenMV/串口监视器输出的传感器初始化成功日志或I\(^2\)C扫描结果（若有）。图7-4建议为“TFT本地显示界面照片”，要求屏幕上能看到温度、湿度、气压、VOC与网络状态字段。
\begin{figure}[H]
\centering
\fbox{\parbox{0.92\linewidth}{\centering 图7-2占位：硬件连接与实物连线图。\\
建议包含：逻辑连接图（模块+接口+引脚号）与实物照片（实际连线），标注UART/I2C/SPI与电源走向。}}
\caption{硬件连接与实物连线图（建议配图）}
\label{fig:hardware_wiring}
\end{figure}
\begin{figure}[H]
\centering
\fbox{\parbox{0.92\linewidth}{\centering 图7-3占位：采集侧传感器调试与在线性验证。\\
建议给出串口日志/调试截图，体现AHT20/BMP280/SGP40初始化成功及读数稳定输出。}}
\caption{采集侧传感器调试与在线性验证（建议配图）}
\label{fig:sensor_debug}
\end{figure}
\begin{figure}[H]
\centering
\fbox{\parbox{0.92\linewidth}{\centering 图7-4占位：网关TFT本地显示界面实拍。\\
要求画面能看清温度、湿度、气压、VOC、气体传感器值与WiFi连接状态等字段。}}
\caption{网关TFT本地显示界面（建议配图）}
\label{fig:tft_ui}
\end{figure}

\subsection{软件实现与调试}
\par 软件实现以端到端链路闭环为目标，重点完成“采集—上报—存储—展示—控制—响应”的全流程打通。嵌入式侧采用模块化类封装方式组织代码，采集侧负责传感器初始化与周期采样，网关侧负责串口解析、MQ气体采样、显示刷新、HTTP上报与MQTT控制。OpenMV侧采用命令驱动脚本，在空闲、采集与识别三种模式之间切换，并按命令输出文本响应或JPEG图像数据。服务器侧采用Express提供数据接收与查询接口，并通过MQTT管理器提供“发送命令并等待响应”的能力，同时将识别结果等异步消息持久化为可查询记录。前端页面通过携带token的请求调用API，完成仪表板绘图、设置保存与人脸系统交互。
\par 联调阶段重点验证三条关键链路。第一条是“STM32→ESP32”的串口数据链路，验证按行上报与超时丢弃策略在噪声或半包情况下仍可正常恢复。第二条是“ESP32→服务器”的HTTP链路，验证数据字段与时间戳格式正确，服务器能写入并在网页端正常显示。第三条是“服务器→MQTT→ESP32→OpenMV→MQTT→服务器”的控制链路，验证命令白名单校验、超时处理、进度消息连续回传与图像Base64回传等功能。对于图像数据，还需验证MQTT缓冲区限制下的分块传输与前端渲染正确性。
\par 配图建议如下。图7-5建议为“端到端数据链路调试截图”，可以用两列：左列是ESP32串口监视器输出（显示WiFi连接成功、HTTP上传成功日志），右列是服务器控制台输出（显示收到数据的JSON内容与时间戳）。图7-6建议为“Web仪表板页面截图”，要求包含多参数曲线、统计摘要与最后更新时间。图7-7建议为“MQTT控制链路时序示意图”，画法为时序图：Web/Server→MQTT Broker→ESP32→OpenMV→ESP32→Broker→Server/Web，标注命令主题与响应主题以及超时窗口。图7-8建议为“人脸系统页面截图”，要求包含实时图像区域、识别结果列表、控制按钮与当前状态提示。
\begin{figure}[H]
\centering
\fbox{\parbox{0.92\linewidth}{\centering 图7-5占位：端到端数据链路调试截图。\\
建议同时展示ESP32串口监视器日志（WiFi/HTTP上传）与服务器控制台日志（/data接收与写入）。}}
\caption{端到端数据链路调试截图（建议配图）}
\label{fig:e2e_data_debug}
\end{figure}
\begin{figure}[H]
\centering
\fbox{\parbox{0.92\linewidth}{\centering 图7-6占位：Web仪表板页面截图。\\
要求包含多参数曲线（Chart.js）、统计摘要、连接状态与最后更新时间等元素。}}
\caption{Web仪表板页面截图（建议配图）}
\label{fig:web_dashboard}
\end{figure}
\begin{figure}[H]
\centering
\fbox{\parbox{0.92\linewidth}{\centering 图7-7占位：MQTT控制链路时序示意图。\\
建议画成时序图，标注topic（smarthome/command、smarthome/response）与GET\_IMAGE分块回传流程。}}
\caption{MQTT控制链路时序示意图（建议配图）}
\label{fig:mqtt_sequence}
\end{figure}
\begin{figure}[H]
\centering
\fbox{\parbox{0.92\linewidth}{\centering 图7-8占位：人脸系统页面截图。\\
要求包含实时图像、识别框/提示、识别结果列表、采集/识别/停止/获取图像/名单管理按钮与状态栏。}}
\caption{人脸系统页面截图（建议配图）}
\label{fig:face_page}
\end{figure}

\subsection{系统测试方案}
\par 系统测试以验证“功能完整性、性能可达性与运行稳定性”为目标，测试对象覆盖采集侧、网关侧、视觉侧与服务器/前端/语音助手等模块。测试环境包含硬件环境（STM32采集板、ESP32-S3网关、OpenMV、TFT屏与传感器模块）、软件环境（Node.js服务器、浏览器端页面、Python语音助手运行环境）与网络环境（家庭路由器或热点、MQTT Broker可达性）。测试过程中需记录关键日志与时间戳，以便定位链路瓶颈。
\par 测试工具方面，基础测试可使用串口监视器、日志输出与浏览器开发者工具完成；若具备条件，可使用逻辑分析仪或示波器辅助验证UART波特率与数据帧完整性，使用网络抓包工具验证HTTP请求与MQTT消息的时序与载荷。对于人脸识别与图像回传，可通过页面显示结果与服务器记录对比验证一致性。对于语音助手，可通过控制台日志记录ASR转写文本、工具调用记录与TTS播放状态。
\par 用例设计方面，功能测试覆盖环境采集、上报、展示、控制、人脸采集/识别/名单管理/图像获取与语音查询等核心路径；性能测试关注端到端时延、链路响应时间与资源占用；稳定性测试关注长时间运行、网络断开重连、模块重启恢复与异常数据处理。建议为每类测试建立用例表格并记录“输入、步骤、期望输出、实际结果与结论”。
\par 配图建议如下。图7-9建议为“测试环境与工具照片/截图拼图”，包含硬件实物搭建照片、服务器运行窗口、浏览器页面、串口监视器窗口等，形成一张整体测试环境展示图。
\begin{figure}[H]
\centering
\fbox{\parbox{0.92\linewidth}{\centering 图7-9占位：测试环境与工具展示图。\\
建议拼图方式展示：硬件搭建实拍、串口监视器、服务器日志窗口、网页端页面与语音助手控制台。}}
\caption{测试环境与工具展示图（建议配图）}
\label{fig:test_env}
\end{figure}

\subsection{功能测试}
\par 环境监测功能测试主要验证采集侧读数合理性、网关解析正确性与平台展示一致性。测试方法是设置稳定环境或进行可控扰动（如用手捂住温湿度传感器、改变环境通风），观察STM32串口输出、ESP32解析结果、服务器存储数据与网页曲线趋势是否一致，确认字段单位与小数位处理正确。对于VOC与气体传感器，可在安全前提下使用酒精棉片/空气清新剂等产生可观测变化，验证趋势响应与恢复过程。
\par 人脸识别功能测试主要验证采集、识别与名单管理闭环。测试步骤包括：在页面发起采集命令并观察进度消息是否连续更新，采集完成后使用名单列举确认样本数量与姓名目录创建正确；进入识别模式后，观察识别结果是否在页面持续刷新，并验证识别与未知人脸的区分；执行删除命令后再次列举确认数据清除。图像获取功能需验证能在不同模式下返回可显示的JPEG图像，并在网页端正确渲染。
\par 语音交互功能测试主要验证“ASR→LLM→工具→TTS”的闭环。测试中应覆盖环境数据查询、统计摘要查询与天气查询等典型问题，观察ASR最终文本是否正确、LLM是否触发工具调用、工具是否成功返回数据、TTS是否完成播报且无明显卡顿。可通过日志对齐“用户文本—工具调用—接口响应—播报文本”四者一致性，验证语音助手的可控性与正确性。
\par 服务器交互功能测试主要验证HTTP上报、鉴权与MQTT命令通路。测试包括：在无token情况下访问受保护页面是否重定向到登录页；登录后是否可正常拉取数据与保存设置；下发OpenMV命令是否返回响应或超时错误；识别结果是否在服务器侧被记录并可被前端查询展示。

\subsection{性能测试}
\par 性能测试重点关注端到端时延与交互响应。数据采集延迟可通过对齐多个时间戳进行估计：STM32采样周期与上报间隔、ESP32上传周期以及服务器写入时间戳。实际操作中可在网页端观察最新数据刷新时间与本地串口输出节拍的对应关系，验证系统具备秒级更新能力。通信响应时间测试可分别测量UART链路（从上报到解析更新）与MQTT链路（从前端下发命令到收到响应）的大致耗时，并记录在不同网络状态下的变化。
\par 人脸识别效果评估可采用小规模样本统计方式进行描述，例如在不同光照与距离条件下对同一人进行多次识别，统计识别成功次数、误识别次数与未知判定次数，并结合LBP距离阈值讨论误识别原因。若不方便给出精确数值，可给出测试方法、评价指标定义与典型现象分析，保证论文表述严谨。
\par 服务器并发处理能力在原型系统中以单设备为主，但可以通过浏览器多标签页并行刷新与多次请求接口的方式验证服务器在一定并发下仍可正常响应。若需要展示并发能力，可在测试中记录服务器CPU占用与响应时间的变化趋势，作为后续扩展的参考。

\subsection{稳定性测试}
\par 稳定性测试以长时间运行与异常场景恢复为核心。长时间运行测试建议至少持续24小时，观察服务器数据文件是否持续写入且无明显异常增长，网页端是否能持续刷新，网关端是否发生死机或卡顿。对于网关与语音助手等存在队列缓冲的模块，应关注是否存在内存占用持续增长或队列堆积导致的性能下降。
\par 异常场景测试建议覆盖网络断开与重连、MQTT Broker不可达、服务器重启、网关重启与传感器临时失联等情况。重点观察系统是否能自动恢复到可用状态，例如WiFi断开后网关是否按周期重连，MQTT断开后是否自动重连并恢复订阅，服务器重启后是否能重新加载配置并继续接收数据。对于传感器失联或读数异常，系统应能通过缺失值标记与容错逻辑避免异常扩散到全局。
\par 边界条件测试在家庭场景可用“网络抖动与高频操作”替代极端环境，例如在识别模式下频繁触发获取图像、同时刷新页面并修改设置，观察系统是否出现响应超时或卡死，并记录出现问题时的日志用于分析。

\subsection{测试结果分析}
\par 测试结果分析建议从“功能、性能、稳定性”三个维度归纳。功能层面可总结环境采集上报、网页可视化、人脸采集识别与语音查询等核心路径均可闭环运行；性能层面可总结系统具备秒级刷新能力、控制链路具备可感知实时性，并说明图像回传在MQTT缓冲区限制下采用Base64与分块策略后可稳定显示；稳定性层面可总结系统能够在网络波动与重连场景下恢复工作，并说明长时间运行中数据裁剪与队列背压机制对资源控制的作用。
\par 问题整改部分可结合实际联调经历进行描述。例如在早期实现中可能出现串口半包导致解析异常、图像数据过大导致MQTT发布失败、识别结果异步消息与GET\_IMAGE响应混淆等问题，后续通过引入串口超时清空、MQTT分块发布、识别结果过滤与超时策略优化等手段进行修复。该部分可显著提升论文的工程可信度。
\par 配图建议如下。图7-10建议为“典型测试结果展示图”，可以是两张关键截图或曲线：一张展示网页端曲线随时间变化的稳定趋势（截图或导出），另一张展示人脸识别结果列表随时间更新的截图，并在图注中说明测试持续时间与场景条件。
\begin{figure}[H]
\centering
\fbox{\parbox{0.92\linewidth}{\centering 图7-10占位：典型测试结果展示图。\\
建议包含：网页端曲线稳定刷新截图 + 人脸识别结果列表更新截图（可拼图），并标注测试持续时间与场景。}}
\caption{典型测试结果展示图（建议配图）}
\label{fig:test_results}
\end{figure}
